\Article{54}
Es sei $h$ eine primitive Wurzel für den Modul $m$, mit der Norm $=p$. Die Glieder der Folge
\[
1,h,hh,h^3\dotsfour h^{p-2}
\]
sind dann einander inkongruent. Daraus folgt leicht, dass jede durch den Modul nicht teilbare ganze Zahl einem dieser Glieder kongruent sein muss; diese Folge stellt also ein vollständiges System inkongruenter Reste unter Ausschluss der Null dar. Der Exponent der Potenz, der eine gegebene Zahl kongruent ist, kann ihr \emph{Index} heißen, wobei $h$ als \emph{Basis} betrachtet wird. Hier folgen einige Beispiele, in denen wir neben jeden Index den absolut kleinsten Rest gesetzt haben.

\begin{center}\emph{Erstes Beispiel.}\end{center}
\[
m=5+4i,\quad p=41,\quad h=1+2i
\]
\begingroup\small
\[
\setlength{\arraycolsep}{2pt}
\begin{array}{r|r||r|r||r|r||r|r||r|r}
\text{Ind.}&\text{Rest}&\text{Ind.}&\text{Rest}&\text{Ind.}&\text{Rest}&\text{Ind.}&\text{Rest}&\text{Ind.}&\text{Rest}\\\hline
0&+1&8&-4&16&-2+2i&24&+2i&32&+1+i\\
1&+1+2i&9&-3+i&17&-1+2i&25&-3i&33&+1+3i\\
2&+1-i&10&-i&18&+4i&26&+2+2i&34&+2\\
3&+3+i&11&+2-i&19&+1+3i&27&+2+i&35&-3\\
4&-2i&12&-1-i&20&-1&28&+4&36&+2-2i\\
5&+3i&13&+1-3i&21&-1-2i&29&+3-i&37&+1-2i\\
6&-2-2i&14&-2&22&-1+i&30&+i&38&-4i\\
7&-2-i&15&+3&23&-3-i&31&-2+i&39&-1-3i
\end{array}
\]
\endgroup

\SourceLeaf{135}%
\begin{center}\emph{Zweites Beispiel.}\end{center}
\[
m=7,\quad p=49,\quad h=1+2i
\]
\begingroup\small
\[
\setlength{\arraycolsep}{2pt}
\begin{array}{r|r||r|r||r|r||r|r||r|r}
\text{Ind.}&\text{Rest}&\text{Ind.}&\text{Rest}&\text{Ind.}&\text{Rest}&\text{Ind.}&\text{Rest}&\text{Ind.}&\text{Rest}\\\hline
0&+1&10&-1-i&20&+2i&30&+2-2i&40&+3\\
1&+1+2i&11&+1-3i&21&+3+2i&31&-1+2i&41&+3-i\\
2&-3-3i&12&-i&22&-1+i&32&+2&42&-2-2i\\
3&+3-2i&13&+2-i&23&-3-i&33&+2-3i&43&+2+i\\
4&-3i&14&-3+3i&24&-1&34&+1+i&44&-2i\\
5&-1-3i&15&-2-3i&25&-1-2i&35&-1+3i&45&-3-2i\\
6&-2+2i&16&-3&26&+3+3i&36&+i&46&+1-i\\
7&+1-2i&17&-3+i&27&-3+2i&37&-2+i&47&+3+i\\
8&-2&18&+2+2i&28&+3i&38&+3-3i&&\\
9&-2+3i&19&-2-i&29&+1+3i&39&+2+3i&&
\end{array}
\]
\endgroup

\Article{55}
Wir fügen einige Bemerkungen über primitive Wurzeln und die Berechnung von Indizes hinzu und lassen die Beweise ihrer Einfachheit wegen weg.

I. In einem gegebenen System entsprechen nach dem Modul $p-1$ kongruente Indizes nach dem Modul $m$ kongruenten Resten, und umgekehrt.

II. Reste, die zu $p-1$ teilerfremden Indizes entsprechen, sind selbst primitive Wurzeln, und umgekehrt.

III. Wird eine primitive Wurzel $h$ als Basis genommen und ist der Index einer anderen primitiven Wurzel $h'$ gleich $t$ und umgekehrt $t'$ der Index von $h$, wenn $h'$ als Basis genommen wird, so gilt $tt'\equiv1\pmod{p-1}$. Sind unter denselben Voraussetzungen die Indizes einer anderen Zahl in diesen beiden Systemen der Reihe nach $u,u'$, so gilt $tu'\equiv u,t'u\equiv u'\pmod{p-1}$.

IV. Werden die Zahlen $1,1+i$ und jeweils ihre drei Assoziierten als zu trivial von den betrachteten Moduln ausgeschlossen, so bleiben die Primzahlen übrig, die wir in Artikel~34 an dritter und vierter Stelle aufgeführt haben. Die Normen der letzteren sind reelle Primzahlen der Form $4n+1$, die der ersteren dagegen Quadrate ungerader reeller Primzahlen; in beiden Fällen ist also $p-1$ durch $4$ teilbar.

V. Bezeichnen wir den Index der Zahl $-1$ mit $u$, so gilt $2u\equiv0\pmod{p-1}$, also entweder $u\equiv0$ oder $u\equiv\frac12(p-1)$. Da aber der Index $0$ dem Rest $+1$ entspricht, muss der Index der Zahl $-1$ notwendig $\frac12(p-1)$ sein.

VI. Bezeichnen wir ebenso mit $u$ den Index der Zahl $i$, so gilt $2u\equiv\frac12(p-1)\pmod{p-1}$, also entweder $u\equiv\frac14(p-1)$ oder $u\equiv\frac34(p-1)$. Hier hängt die Mehrdeutigkeit jedoch von der Wahl der primitiven Wurzel ab. Ist nämlich bei der als Basis genommenen primitiven Wurzel $h$
\SourceLeaf{136}%
der Index der Zahl $i$ gleich $\frac14(p-1)$, so wird er $\frac34(p-1)$, wenn $h^\mu$ als Basis genommen wird, wobei $\mu$ eine positive ganze Zahl der Form $4n+3$ bezeichnet, die zu $p-1$ teilerfremd ist, zum Beispiel die Zahl $p-2$ selbst; und umgekehrt. Daher gibt die eine Hälfte der primitiven Wurzeln der Zahl $i$ den Index $\frac14(p-1)$, die andere den Index $\frac34(p-1)$. Für jene Basen hat offenbar $-i$ den Index $\frac34(p-1)$, für diese den Index $\frac14(p-1)$.

VII. Ist der Modul eine positive reelle Primzahl der Form $4n+3$, etwa $=q$, und somit $p=qq$, so sind die Indizes aller reellen Zahlen durch $q+1$ teilbar. Bezeichnet nämlich $t$ den Index der reellen Zahl $k$, so gilt wegen $k^{q-1}\equiv\pmod q$ auch $(q-1)t\equiv0\pmod{qq-1}$, und somit ist $\frac{t}{q+1}$ eine ganze Zahl. Ebenso sind die Indizes rein imaginärer Zahlen wie $ki$ durch $\frac12(q+1)$ teilbar. Daher ist klar, dass primitive Wurzeln für solche Moduln nur unter gemischten Zahlen zu suchen sind.

\Needspace{8\baselineskip}VIII. Dagegen lassen sich für einen Modul $m$, der eine gemischte komplexe Primzahl ist und dessen Norm $p$ folglich eine reelle Primzahl der Form \mbox{$4n+1$} ist, beliebige primitive Wurzeln auch unter reellen Zahlen wählen, unter denen sich sogar ein vollständiges System inkongruenter Reste angeben lässt (Artikel~40). Offenbar ist jede reelle Zahl, die für den komplexen Modul $m$ eine primitive Wurzel ist, zugleich in der Arithmetik reeller Zahlen eine primitive Wurzel für den Modul $p$, und umgekehrt.

